\setcounter{page}{28}
\section{\S3. Localization of Categories}

\textbf{Definition.}
Let \(\cat{C}\) be a category. A collection \(S\) of morphisms of \(\cat{C}\) is called a \textit{multiplicative system} if it satisfies the following axioms \(\text{FR1}\)--\(\text{FR3}\):

\begin{enumerate}
\item[(FR1)] If \(f,g\in S\) and the composite \(fg\) exists, then \(fg\in S\). For any \(X\in\operatorname{Ob}\cat{C}\), \(\id_X\in S\).

\item[(FR2)] Any diagram
\[
X\xrightarrow{u} Y \xleftarrow{s} Z,\quad s\in S
\]
can be completed to a commutative diagram
\[
\begin{CD}
W @>v>> Z\\
@VtVV @VsVV\\
X @>u>> Y
\end{CD}
\]
with \(t\in S\). The analogous statement holds with all arrows reversed.

\item[(FR3)] Let \(f,g\colon X\to Y\) be morphisms in \(\cat{C}\). The following are equivalent:
\begin{enumerate}
\item[(i)] There exists \(s\colon Y\to Y'\) in \(S\) such that \(sf = sg\);
\item[(ii)] There exists \(t\colon X'\to X\) in \(S\) such that \(ft = gt\).
\end{enumerate}
\end{enumerate}

\setcounter{page}{29}
\textbf{Definition.}
Let \(\cat{C}\) be a category and \(S\) a multiplicative system of morphisms in \(\cat{C}\). The \textit{localization of \(\cat{C}\) with respect to \(S\)} is a category \(\cat{C}_S\) together with a functor \(Q\colon\cat{C}\to\cat{C}_S\) such that:
\begin{enumerate}
\item[(a)] \(Q(s)\) is an isomorphism for every \(s\in S\);
\item[(b)] Any functor \(F\colon\cat{C}\to\cat{D}\) sending each \(s\in S\) to an isomorphism factors uniquely through \(Q\).
\end{enumerate}

\textbf{Remark.} One can prove such a localization exists without extra assumptions on \(S\), though we do not use this fact here.

\medskip
\textbf{Proposition 3.1.}
Let \(\cat{C}\) be a category and \(S\) a multiplicative system in \(\cat{C}\). The localization \(\cat{C}_S\) may be constructed as follows:
\(\operatorname{Ob}\cat{C}_S = \operatorname{Ob}\cat{C}\), and for all \(X,Y\in\operatorname{Ob}\cat{C}\),
\[
\Hom_{\cat{C}_S}(X,Y) = \varinjlim_{I_X} \Hom_{\cat{C}}(X',Y),
\]
where \(I_X\) is the category whose objects are morphisms \(s\colon X'\to X\) in \(S\), and morphisms are commutative triangles over \(X\).

\setcounter{page}{30}
Furthermore, if \(\cat{C}\) is additive, then \(\cat{C}_S\) is also additive.

\textit{Proof.}
Using \(\text{FR1}\), \(\text{FR2}\), \(\text{FR3}\), the index category \(I_X\) satisfies the inductive limit axioms (LI, L2, L3) of [GT, Ch.\ I, \S], so the direct limit is well-behaved.

A morphism \(X\to Y\) in \(\cat{C}_S\) is represented by a right fraction \(X\xleftarrow{s}X'\xrightarrow{f}Y\) with \(s\in S\). Two such fractions define the same morphism iff there exist \(u\colon X'''\to X'\), \(g\colon X'''\to X''\) in \(S\) making the obvious diagram commute and equalizing the representatives.

\setcounter{page}{31}
Composition of fractions is defined via \(\text{FR2}\): given fractions \(X\xleftarrow{s}X'\xrightarrow{a}Y\) and \(Y\xleftarrow{t}Y'\xrightarrow{b}Z\), complete the diagram to get a common lift and form the composite morphism \(X\to Z\). The result is independent of representative and lift.

The localization functor \(Q\colon\cat{C}\to\cat{C}_S\) satisfies the universal property, and additivity of \(\cat{C}\) passes to \(\cat{C}_S\) since the limit preserves abelian group structure on Hom-sets.

\medskip
\textbf{Definition.}
Let \(\cat{C}\) be a triangulated category and \(S\) a multiplicative system in \(\cat{C}\). We say \(S\) is \textit{compatible with the triangulation} if:
\begin{enumerate}
\item[(FR4)] \(s\in S \implies T(s)\in S\), where \(T\) is the translation functor;
\item[(FR5)] Axiom \(\text{TR3}\) holds with the extra condition: if \(f,g\in S\) are the given morphisms of triangles, then one may choose the filler \(h\in S\).
\end{enumerate}

\setcounter{page}{32}
\textbf{Proposition 3.2.}
If \(\cat{C}\) is triangulated and \(S\) a triangulation-compatible multiplicative system, then \(\cat{C}_S\) admits a \textit{unique} triangulated structure such that \(Q\colon\cat{C}\to\cat{C}_S\) is a triangle functor, and \(Q\) satisfies the universal property for triangle functors into other triangulated categories.

\textit{Proof.} Left to the reader.

One may also compute morphism sets as
\[
\Hom_{\cat{C}_S}(X,Y) = \varinjlim_{J_Y}\Hom_{\cat{C}}(X,Y'),
\]
where \(J_Y\) has objects \(s\colon Y\to Y'\in S\) and morphisms commutative diagrams under \(Y\).

\medskip
\textbf{Proposition 3.3.}
Let \(\cat{C}\) be a category, \(S\) a multiplicative system in \(\cat{C}\), and \(\cat{D}\subseteq\cat{C}\) a full subcategory such that \(S\cap\cat{D}\) is a multiplicative system in \(\cat{D}\). Assume one of the following:
\begin{enumerate}
\item[(i)] For every \(s\colon X'\to X\) in \(S\) with \(X\in\operatorname{Ob}\cat{D}\), there exists \(f\colon X''\to X'\) with \(X''\in\operatorname{Ob}\cat{D}\) and \(sf\in S\);
\item[(ii)] The same with all arrows reversed.
\end{enumerate}
Then the natural functor
\[
\cat{D}_{S\cap\cat{D}} \to \cat{C}_S
\]
is fully faithful, so \(\cat{D}_{S\cap\cat{D}}\) may be identified with a full subcategory of \(\cat{C}_S\).

\setcounter{page}{33}
\textit{Proof.} Straightforward.

\medskip
\textbf{Proposition 3.4.}
Let \(\cat{C}\) be a category, \(S\) a multiplicative system, \(Q\colon\cat{C}\to\cat{C}_S\) the localization functor. Let \(\cat{D}\) be any category, and \(F,G\colon\cat{C}_S\to\cat{D}\) functors. Then the natural restriction map
\[
\alpha\colon \Hom(F,G) \to \Hom(FQ,GQ)
\]
between functor morphism sets is bijective.

\setcounter{page}{34}
\textit{Proof.}
A functor morphism \(F\to G\) assigns to each object \(X\in\cat{C}_S\) a morphism \(F(X)\to G(X)\) making all morphism squares commute.

Injectivity of \(\alpha\) is clear since \(\operatorname{Ob}\cat{C}_S=\operatorname{Ob}\cat{C}\). For surjectivity, take a morphism \(FQ\to GQ\). Since every morphism in \(\cat{C}_S\) is a fraction built from morphisms of \(\cat{C}\) and elements of \(S\), and \(F(s),G(s)\) are isomorphisms for \(s\in S\), the given functor morphism on \(FQ,GQ\) extends uniquely to a functor morphism on \(F,G\) over \(\cat{C}_S\).